This paper is about implementing \emph{functional logic programming} (FLP)
\parencite{antoy_2010}. FLP extends functional programming with nondeterministic
choice, logical variables, and equational constraints. These features allow a
declarative style of programming in which the programmer describes relationships
between values, and the language searches for solutions. For example, the last
element of a list can be computed by searching for a decomposition:

\begin{haskellcode}
last :: [Nat] -> Nat
last ys = exists xs :: [Nat]. exists x :: Nat. cat xs [x] =:= ys. x.
\end{haskellcode}

\noindent This program introduces logical variables \mintinline{haskell}|xs| and
\mintinline{haskell}|x|, constrains them so that concatenating
\mintinline{haskell}|xs| with the singleton \mintinline{haskell}|[x]| yields
\mintinline{haskell}|ys|, and returns \mintinline{haskell}|x|. The language
discovers the appropriate values automatically through \emph{narrowing}: a
progressive refinement of logical variables driven by the demands of the
computation.

In a companion paper \parencite{jones_2026} we studied the semantics of FLP
through the lens of \emph{algebraic effects} \parencite{plotkin_2003}. The key
observation of that work is that \textbf{the characteristic constructs of FLP
are effects} in the sense of \citeauthor{plotkin_2003}. This led to a core
calculus based on Call-by-Push-Value (CBPV) \parencite{levy_2003}, a
domain-theoretic denotational semantics based on the lower powerdomain, and a
comprehensive equational theory that validates a range of program
transformations. Crucially, the denotational semantics gives rise---via a
\emph{representation theorem}---to an abstract machine whose transitions are
justified by the semantics. The equational theory, in turn, provides a catalogue
of valid program equations.

The present paper asks: \emph{can this theory serve as a practical guide for
implementation?} We answer in the affirmative by presenting
\Gweek{},\!\footnote{%
  Available at \url{https://github.com/xxx/gweek}. The name is a portmanteau of
  the Welsh \emph{gweek} (`search') and the characteristic exclamation of the
  language user upon discovering a particularly elegant solution.}
a small but complete FLP language whose implementation is derived directly from
the theory. \Gweek{} implements a variant of the abstract machine of
\parencite{jones_2026}, and because every transition of the machine is validated
by the denotational semantics, the machine is \emph{evidently correct}: its
correctness is not established by ad hoc reasoning, but follows directly from
the semantic model.

Three aspects of this connection are particularly illuminating.

First, the equational theory of \parencite{jones_2026} serves directly as the
rewrite rules of a \emph{peephole optimiser}
(\cref{section:optimisation}). Each rewrite rule applied by the optimiser is a
validated equation of the theory; its soundness follows from the full abstraction
theorem of the denotational semantics.

Second, the CBPV distinction between values and computations leads naturally to a
\emph{suspension mechanism} for \texttt{let} bindings
(\cref{section:suspensions}). This mechanism defers the evaluation of bound
computations, dramatically improving termination by allowing constraints to prune
the search tree before the full nondeterministic tree is materialised. The
equational theory illuminates precisely when and why this matters.

Third, the machine's treatment of logical variables and narrowing
(\cref{section:machine}) is a direct realisation of the abstract machine of
\parencite{jones_2026}. A CBPV term gives rise to a nondeterministic search
tree; the implementation explores this tree using one of four strategies. The
narrowing strategy---which refines a logical variable only when a case expression
demands its value---corresponds to \emph{weak head narrowing}, and inherits the
completeness guarantee of \parencite{jones_2026}.

\paragraph{Contributions}

The contributions of this paper are:
\begin{itemize}
  \item A description of \Gweek{}, a typed higher-order FLP language. We present
    its surface syntax, a CBV core calculus, and its translation to CBPV
    (\cref{section:source,section:ir}).
  \item A presentation of the equational theory of CBPV with FLP effects, its
    restriction to CBV equations, and an example proof showing how CBPV equations
    validate CBV transformations (\cref{section:ir}).
  \item A CEK-style abstract machine for CBPV that explores nondeterministic
    search trees, with four search strategies (\cref{section:machine}).
  \item A description of the Rust implementation, including arena allocation,
    persistent data structures, and the compilation pipeline
    (\cref{section:implementation}).
  \item A peephole optimiser whose rewrite rules are drawn from the equational
    theory (\cref{section:optimisation}).
  \item A suspension mechanism that improves termination by deferring
    evaluation, with an analysis of its interaction with nondeterminism
    (\cref{section:suspensions}).
\end{itemize}
